% Eksempel på kapittelinnhold.
% Bytt ut med ditt eget.

\chapteropening{D}{ette er åpningssetningen}
til kapittelet, satt med initial og kapiteler på de første
ordene. Åpningsavsnittet bør være langt nok til å fylle ut
den vertikale plassen initialen reserverer, ellers vil neste
seksjonsoverskrift krasje med initialen.

Et nytt avsnitt fortsetter uten ekstra avstand. Klassisk
boktypografi bruker innrykkede avsnitt og minimalt med
vertikalt mellomrom --- det er prosaens rytme som drar
leseren gjennom.


\section{Første seksjon}

En seksjonsoverskrift uten nummer. Bokmalen bruker
unummererte seksjoner slik at kapittel- og delnummerene
bærer det strukturelle hierarkiet alene. Seksjonene vises
fortsatt i innholdsfortegnelsen, så navigasjon er bevart.

Siteringer i denne malen vises som små superskript-tall,
plassert nøyaktig der de hører hjemme i teksten\cite{example-article}.
Flere kilder kan grupperes innenfor én sitering og blir
sortert og komprimert automatisk\cite{example-article,example-book}.


\section{Matematikk}

Matematikk settes inline, $E = mc^2$, eller som en
nummerert ligning:
\begin{equation}
  i\hbar \frac{\partial \psi}{\partial t}
  = \hat{H}\psi.
  \label{eq:schrodinger}
\end{equation}
Ligning \eqref{eq:schrodinger} er den tidsavhengige
Schr\"odingerligningen, skrevet her kun som et eksempel
på en merket ligning som kan refereres fra andre steder
i kapittelet.

En andre ligning kan følge med en annen etikett:
\begin{equation}
  \langle A \rangle_\psi = \langle \psi | \hat{A} | \psi \rangle.
  \label{eq:forventning}
\end{equation}
Avstanden rundt ligninger er satt i \texttt{bokstil.sty}
og gir matematikken nok luft til å puste uten å miste
kontakt med teksten omkring.

\subsection{En underseksjon}

Underseksjoner nestes under seksjoner og er også
unummererte. De er nyttige når en seksjon blir lang nok
til å trenge intern struktur, men ikke lang nok til å
fortjene sin egen seksjon.


\section{Figurer og tabeller}

En figur inkluderes ved å plassere filen i \texttt{figurer/}-mappen
og kalle \texttt{\textbackslash includegraphics}:

\begin{figure}[t]
  \centering
  \includegraphics[width=0.7\linewidth]{Energy_levels}
  \caption{Et eksempel på en figur som viser energinivåer.}
  \label{fig:nivaer}
\end{figure}

Henvis til figurer med figur~\ref{fig:nivaer}. Plasseringen
\texttt{[t]} lar figuren stige til toppen av en side, som er
den vanlige bokkonvensjonen --- figurer sitter sjelden
inline i langformet prosa.

En tabell kan oppsummere nøkkelverdier:
\begin{table}[t]
  \centering
  \caption{Fundamentale fysiske konstanter.}
  \label{tab:konstanter}
  \begin{tabular}{lcr}
    \toprule
    Størrelse           & Symbol  & Verdi                         \\
    \midrule
    Lysets hastighet    & $c$     & $2{,}998 \times 10^8$ m/s     \\
    Plancks konstant    & $\hbar$ & $1{,}055 \times 10^{-34}$ Js  \\
    Elementærladning    & $e$     & $1{,}602 \times 10^{-19}$ C   \\
    \bottomrule
  \end{tabular}
\end{table}

% For native TikZ-figurer i en bok, foretrekk standalone-arbeidsflyten:
% kompiler hver figur i sin egen fil med \documentclass{standalone},
% og inkluder den resulterende PDF-en med \includegraphics.
% Dette holder kapittelfiler rene og unngår treg rekompilering av
% hele boken når en figur endres. bokstil.sty definerer ingen
% figurpalett; vil du bruke samlingens farger fra notat- og
% handout-stilene, definer dem i standalone-preamblen med xcolor:
%   \definecolor{darkorange}{RGB}{184, 92,  0}
%   \definecolor{darkolive} {RGB}{ 74,107, 18}

Tabeller i denne malen bruker \texttt{booktabs}-konvensjonen:
kun horisontale streker, brukt sparsomt --- topp, midt og
bunn. Vertikale streker er bevisst utelatt.


\section{Avslutning}

En avsluttende seksjon kan runde av kapittelet og peke
mot stoffet i det neste. I en lengre bok er dette stedet
hvor man kan forhåndsvise neste kapittel eller ta et
skritt tilbake for å oppsummere det som er dekket.
